\subsubsection{Кусково-поліноміальна інтерполяція: локальний барицентричний метод}

Глобальна барицентрична інтерполяція, розглянута у попередніх розділах, використовує значення всіх пікселів зображення для обчислення кольору кожної нової точки. З математичної точки зору це означає побудову інтерполяційного полінома надзвичайно високого ступеня (наприклад, ступеня 999 для рядка з 1000 пікселів). Хоча використання вузлів Чебишева нівелює феномен Рунге, на зображеннях із різкими просторовими переходами (наприклад, знімки екрана з текстом, графічні інтерфейси, термінали) глобальний підхід може провокувати появу артефактів дзвону (ringing artifacts, ефект Гіббса). Це проявляється у вигляді загасаючих хвилеподібних ліній (ехо-контурів), що поширюються від різких контрастних меж на сусідні області зображення.

Для усунення цієї проблеми було розроблено та програмно реалізовано алгоритм кусково-поліноміальної інтерполяції — локальний барицентричний метод (\texttt{processLocalBarycentric}). Його фундаментальна відмінність полягає у використанні концепції ковзного вікна (sliding window) фіксованого розміру $K \times K$, де $K$ — параметр згладжування (зазвичай 3, 4 або 6). 

Замість глобальної інтерполяції, значення цільового пікселя обчислюється виключно на основі його найближчого локального оточення. Алгоритмічний конвеєр локального методу складається з наступних етапів:
\begin{enumerate}
	\item \textbf{Проекція координат:} Цільові координати нормалізуються в діапазон $[-1, 1]$, після чого проектуються назад на неперервний простір вихідного зображення для знаходження умовної точки $(x_{src}, y_{src})$.
	\item \textbf{Визначення локального вікна:} Навколо знайденої точки обчислюються цілочисельні індекси початкового кута вікна $(x_{start}, y_{start})$. Розмір вікна жорстко обмежується краями вихідного зображення за допомогою функції \texttt{std::clamp} (або \texttt{std::min/max}) для уникнення виходу за межі пам'яті (out-of-bounds memory access).
	\item \textbf{Локальна нормалізація:} Координати умовної точки $(x_{src}, y_{src})$ перераховуються у локальну систему координат вікна $K \times K$, перетворюючись на локальні змінні $local\_x$ та $local\_y$ у діапазоні $[-1, 1]$.
	\item \textbf{Барицентрична оцінка:} Для обчислення фінального кольору застосовується класична двовимірна барицентрична формула \eqref{eq:barycentric_2d_formula}, але межі підсумовування обмежуються розміром вікна (від $0$ до $K-1$), а як значення $f_{i,j}$ використовуються виключно пікселі з поточного вікна.
\end{enumerate}

Локальні просторові вузли для вікна $K \times K$ генеруються за тим самим законом Чебишева-Лобатто \eqref{eq:chebyshev_nodes}, але адаптовані під кількість елементів $K$. Обчислення вагових коефіцієнтів та самої суми відбувається безпосередньо під час проходу по пікселях. Щоб уникнути ділення на нуль, алгоритм перевіряє потрапляння локальних координат $local\_x$ та $local\_y$ у вузли сітки (з точністю до $10^{-11}$). У разі точного збігу, колір відповідного пікселя з вихідного масиву копіюється у цільове зображення.

Використання локального вікна кардинально змінює обчислювальну складність алгоритму. Оскільки обчислення кольору кожної цільової точки вимагає перебору лише $K \times K$ сусідів, складність процесу масштабування стає повністю незалежною від розмірів вихідного зображення і становить $\mathcal{O}(H_2 \cdot W_2 \cdot K^2)$. 

Враховуючи, що $K$ є малою константою, локальний метод демонструє надзвичайно високу швидкість виконання, поступаючись лише вбудованим методам найближчого сусіда або білінійної інтерполяції. Програмна імплементація додатково оптимізована за рахунок багатопотоковості: матриця цільового зображення розбивається на горизонтальні сегменти, кожен з яких паралельно обробляється окремим потоком з пулу \texttt{std::thread}. 

Експериментальні дослідження підтвердили, що локальний барицентричний метод повністю усуває ефект дзвону на різких градієнтах, забезпечуючи ідеальне збереження контрастних меж при масштабуванні технічної графіки та тексту.